\section{Introducción}
La conectividad de largo alcance y bajo consumo que ofrece la tecnología \textbf{LoRa} (Long Range) la convierte en una pieza clave para aplicaciones de Internet de las Cosas (IoT) en entornos de monitoreo remoto. En este laboratorio se integrarán sensores de gas inflamable y de temperatura con un módulo de radio LoRa para diseñar un sistema IoT capaz de detectar gases peligrosos y monitorear variables ambientales a distancia. El sistema enviará inalámbricamente los datos de los sensores mediante LoRa, permitiendo la supervisión en tiempo real de condiciones potencialmente peligrosas a lo largo de varios kilómetros de distancia. Además, incluirá actuadores de alerta (buzzer, LED y servomotor) que se activan cuando los sensores superen umbrales predefinidos, emulando un sistema de alarma de fuga de gas o incendios.

LoRa es una tecnología de modulación por espectro ensanchado que opera en bandas sub-GHz libres y puede lograr comunicaciones de varios kilómetros con un consumo muy bajo de energía. Esto la hace ideal para escenarios sin cobertura WiFi/celular, como la detección temprana de incendios forestales o fugas de gas en instalaciones remotas. De hecho, numerosos estudios han demostrado la eficacia de LoRa en sistemas de detección de gases peligrosos y monitoreo ambiental. Por ejemplo, Sulaiman \textit{et al.} implementaron una red de sensores LoRa con \textbf{MQ-2} (gas combustible) y \textbf{DHT11} (temperatura/humedad) para detección de incendios forestales, logrando transmitir datos de humo y temperatura en tiempo real a más de 1 km sin infraestructura de internet. González \textit{et al.} desarrollaron una red LoRa de bajo costo para monitorear calidad del aire y detectar gases tóxicos, combinando sensores comerciales (p.ej. MQ) con sensores de grafeno para medir compuestos volátiles y NO$_2$ en ambientes urbanos. Asimismo, Herring \textit{et al.} demostraron nodos LoRaWAN subterráneos con termocupla capaces de sobrevivir incendios de maleza y enviar mediciones de temperatura de suelo de forma confiable durante un incendio. Estos trabajos del estado del arte evidencian las ventajas de LoRa para sistemas de alerta temprana: amplia cobertura, operación en áreas sin comunicación tradicional, y larga vida de baterías para despliegues autónomos. En este laboratorio, se aprovecharán estas propiedades de LoRa integrando los componentes descritos a continuación, para construir un prototipo funcional de un sistema IoT de monitoreo de gas y temperatura con alarmas locales y transmisión inalámbrica de datos.

\subsection*{Objetivo general}
Diseñar e implementar un sistema IoT de alerta y monitoreo ambiental que utilice tecnología LoRa para la comunicación de largo alcance entre un nodo sensor y una estación receptora. El nodo sensor integrará un módulo LoRa \textbf{RAK3272S} conectado a un microcontrolador \textbf{Arduino Mega 2560} con sensores de gas inflamable (\textbf{MQ-2}) y temperatura (\textbf{LM35}), así como actuadores de alarma (buzzer, servomotor y LED), con el fin de detectar gases peligrosos y variaciones de temperatura en entornos remotos, activando alertas locales y transmitiendo los datos de forma inalámbrica en tiempo real.

\subsection*{Objetivos específicos}
\begin{itemize}
\item Integración de sensores y actuadores: Montar el circuito del nodo IoT conectando adecuadamente el sensor de gas MQ-2, el sensor de temperatura LM35, un zumbador (\textit{buzzer}), un servomotor y un LED al Arduino Mega, garantizando su correcta alimentación y lectura/señal.
\item Configuración del módulo LoRa (RAK3272S): Configurar el transceptor LoRa RAK3272S para operar en modo LoRaWAN o punto a punto según el escenario del laboratorio, definiendo parámetros de comunicación (frecuencia, SF, potencia) y enlazándolo con el Arduino Mega mediante comunicación serial UART. Se utilizará un conversor \textbf{USB-TTL} para programar o enviar comandos AT al módulo LoRa de ser necesario.
\item Programación del microcontrolador: Desarrollar el código en Arduino C/C++ para leer periódicamente los valores analógicos del MQ-2 y LM35, procesarlos (por ejemplo, calcular ppm estimadas o grados Celsius) y decidir cuándo activar las alarmas (hacer sonar el buzzer, mover el servo, encender el LED) según umbrales de seguridad predefinidos. El código también deberá enviar los datos de sensor y eventos de alarma a través del módulo LoRa.
\item Transmisión y recepción de datos: Establecer comunicación inalámbrica LoRa de largo alcance entre el nodo sensor y un receptor (que puede ser otro módulo LoRa conectado a una PC/servidor o una pasarela LoRaWAN). Comprobar que los datos enviados (lecturas de gas/temperatura, estado de alarmas) se reciben correctamente en el destino, ya sea visualizándolos en un monitor serial o en una plataforma IoT.
\item Validación del sistema de alerta: Realizar pruebas controladas exponiendo el sensor MQ-2 a gas (por ejemplo, gas butano de encendedor) y variando la temperatura del LM35, para verificar que el sistema activa el buzzer, LED y servo cuando las lecturas exceden los umbrales, y que dichas condiciones de alarma son transmitidas exitosamente vía LoRa hacia el receptor remoto.
\end{itemize}


% TODO: Completar 